\chapter{Background}
\label{ch:background}

\section{Neural Network Quantization}

Quantization maps high-precision floating-point tensors to lower-precision integer representations. A real value $x$ can be represented by an integer value $q$ and a scale $s$:

\begin{equation}
  q = \mathrm{clip}\left(\mathrm{round}\left(\frac{x}{s}\right), q_{\min}, q_{\max}\right),
  \qquad
  \hat{x} = q \cdot s.
\end{equation}

This thesis considers fake quantization, int16 reference inference, and selected int8 overrides for sensitivity analysis.

\section{Fixed-Point Arithmetic}

The RTL-oriented flow uses fixed-point formats such as Q8.8, Q0.16, and Q1.15. For Q8.8, a signed 16-bit integer $q$ represents:

\begin{equation}
  x = \frac{q}{2^8}.
\end{equation}

Multiplication changes the number of fractional bits. For example, Q8.8 multiplied by Q8.8 produces Q16.16 and must be requantized before being stored again as Q8.8.

\section{Rounding and Saturation}

Rounding and saturation are important for bit-true behavior. The implementation uses round-to-nearest-even in key requantization paths and clamps values to the target integer range.

\section{Mamba Block Structure}

The evaluated model contains repeated Mamba-style blocks. A block contains normalization, input projection, activation paths, a selective scan state update, element-wise gating, output projection, and residual addition.

\section{Patch Embedding and Projection Layers}

The patch embedding layer projects input patches into the model dimension. Inside each Mamba block, \texttt{in\_proj}, \texttt{dt\_proj}, and \texttt{out\_proj} are projection layers with different sensitivity to low-precision quantization.

\section{Selective Scan and Gate Path}

The selective scan updates a recurrent state:

\begin{equation}
  s_t = \lambda_t s_{t-1} + (1 - \lambda_t) u_t.
\end{equation}

The gate path computes an element-wise product between the selective scan output and the activated $z$ path:

\begin{equation}
  y_{\mathrm{gate}} = \mathrm{SiLU}(z) \odot s.
\end{equation}

\section{Hardware Validation}

The RTL implementation consumes memory initialization files and golden outputs generated by Python. This enables stage-by-stage validation between Python hardware-like simulation and RTL dumps.

